%% Légende du diagramme des cas d'utilisation (uc) : les cinq éléments de base.
%% Bandeau de cinq vignettes séparées par des filets gris : acteur humain, « chose »
%% en interaction, cas d'utilisation, relation «extend», relation «include».
\documentclass[tikz,border=4pt]{standalone}
\usepackage{liaisons}
\begin{document}
\begin{tikzpicture}[x=1cm,y=1cm,
  leg/.style={font=\scriptsize, align=center, text width=1.75cm, anchor=north, inner sep=1pt,
              execute at begin node={\hyphenpenalty=10000\exhyphenpenalty=10000}},
  filet/.style={draw=black!30, line width=0.5pt},
  trait/.style={line width=1pt, line cap=round, line join=round},
  pointille/.style={-{Straight Barb[length=4pt,width=6pt]}, line width=0.9pt,
                    dash pattern=on 3pt off 2pt},
  ster/.style={font=\scriptsize, inner sep=2pt}]

%% ---- macros de pictogrammes -----------------------------------------
%% bonhomme filiforme centré en (#1,#2), couleur #3
\newcommand\acteur[3]{%
  \begin{scope}[shift={(#1,#2)}, trait, draw=#3]
    \draw (0,0.42) circle (0.15);
    \draw (0,0.27) -- (0,-0.12);
    \draw (-0.28,0.12) -- (0.28,0.12);
    \draw (-0.22,-0.5) -- (0,-0.12) -- (0.22,-0.5);
  \end{scope}}
%% petit cube en perspective cavalière centré en (#1,#2), couleur #3
\newcommand\cube[3]{%
  \begin{scope}[shift={(#1,#2)}, trait, draw=#3]
    \draw[fill=#3!12] (-0.34,-0.4) rectangle (0.24,0.18);
    \draw (-0.34,0.18) -- (-0.14,0.4) -- (0.44,0.4) -- (0.24,0.18);
    \draw (0.44,0.4) -- (0.44,-0.18) -- (0.24,-0.4);
  \end{scope}}

%% ---- vignette 1 : acteur humain -------------------------------------
\acteur{0}{1.05}{solideB}
\node[leg] at (0,0.4) {Acteur humain};

%% ---- vignette 2 : « chose » en interaction ---------------------------
\cube{1.85}{1.05}{solideE}
\node[leg] at (1.95,0.4) {« Chose » en interaction (bateau, énergie\dots)};

%% ---- vignette 3 : cas d'utilisation ----------------------------------
\draw[trait, draw=solideD, fill=solideD!12] (3.9,1.05) ellipse (0.95 and 0.42);
\node[font=\scriptsize, align=center, text width=1.6cm] at (3.9,1.05) {Franchir la Loire};
\node[leg] at (3.9,0.4) {Cas d'utilisation};

%% ---- vignette 4 : «extend» -------------------------------------------
\draw[pointille, solideD] (5.85,0.72) -- (5.85,1.42);
\node[ster, anchor=west] at (5.95,1.07) {«extend»};
\node[leg] at (6.05,0.4) {Fonction \textbf{non} indispensable};

%% ---- vignette 5 : «include» ------------------------------------------
\draw[pointille, solideD] (7.9,1.42) -- (7.9,0.72);
\node[ster, anchor=west] at (8.0,1.07) {«include»};
\node[leg] at (8.1,0.4) {Fonction \textbf{indispensable}};

%% ---- filets de séparation --------------------------------------------
\foreach \x in {0.95,2.9,4.95,7.0} { \draw[filet] (\x,-0.55) -- (\x,1.65); }

%% ---- mention finale ---------------------------------------------------
\draw[filet] (-0.95,-0.7) -- (9.2,-0.7);
\node[font=\scriptsize, anchor=north west] at (-0.95,-0.8)
     {acteurs \textbf{principaux} à gauche du système, acteurs \textbf{secondaires} à droite};
\end{tikzpicture}
\end{document}
